\lecture{8}{Tue 30 Sep 2025 12:31}{Expanding Universe}

The simplest metric we can use with an expanding component is
\[
	\mathrm{d} s^2=-\mathrm{d} t^2+a^2(t)\left( \mathrm{d} x^2+\mathrm{d} y^2+\mathrm{d} z^2 \right) 
.\]
This is the simplest case of an FRW universe. Last time we worked through the geodesic equation using the action principle, giving us
\[
	0=\frac{\mathrm{d}}{\mathrm{d}\lambda } \frac{\partial L}{\partial q'} -\frac{\partial L}{\partial q} ,
\]
\[
	q=t,\qquad q=x^i,\qquad q'=\frac{\mathrm{d}q}{\mathrm{d}\lambda } ,\qquad \dot{f} =\frac{\mathrm{d}f}{\mathrm{d}t} 
.\]
\[
	L=\frac{1}{2} \int -g_{\mu \nu } x^{\mu '} x^{\nu '}  \,\mathrm{d} \lambda 
.\]
\[
	\frac{\mathrm{d}}{\mathrm{d}\lambda } a(t)=\dot{a} (t)\frac{\mathrm{d}t}{\mathrm{d}\lambda } 
.\]
We then found
\[
	0=t''+a \dot{a} \left( \frac{t'}{a}  \right) ^2
.\]
\[
	0=\frac{\mathrm{d}}{\mathrm{d}\lambda } \left( a^2(t)x_i' \right) =a^2(t)\left(x_i''+2\frac{\dot{a} (t)}{a(t)} t'x_i'\right)
.\]
Staring at this somehow gives us the Christoffel symbol in $t$
\[
	\Gamma_{\alpha \beta } ^t=a\dot{a} \delta _{\alpha \beta }
.\]
This is because the $t''$ equation is the Geodesic equation
\[
	t''+a \dot{a} \left( \frac{t'}{a}  \right) ^2=q''+\Gamma _{\alpha \beta } ^q x^{\alpha '} x^{\beta '} 
.\]
This is equal probably because they're both 0. We then read it off since they're the same and $t''$ becomes $q''$. For $\Gamma _{\alpha \beta } ^i$ in $x^i$,
\[
	\Gamma _{\alpha \beta } ^i=\frac{\dot{a} }{a}\delta _i^j
.\]

Ok now we want to look at a photon travelling through the expanding universe. Since we cannot use proper time $\tau $, we use an affine parameter $\lambda '=A\lambda +B$. We have a lot of freedom for choosing this affine parameter, in particular the scaling factor $A$. We are going to choose $\lambda $ such that in a fixed reference frame of our choice, $p^\mu =\frac{\mathrm{d}x^\mu }{\mathrm{d}\lambda } $, for $p^\mu $ photon momentum. Such a frame exists. We will always assume the affine parameter for massless particles is chosen in this way.

The energy is just $E=p^0=\frac{\mathrm{d}x^0}{\mathrm{d}\lambda } $. We need to solve the Geodesic equation for this, and then we need to interpret the meaning of this physically. We will do the latter by choosing locally inertial coordinates somewhere else, and then use general covariance.

The energy of a photon has observer dependence (doppler shift). Given a 4-momentum $U^\mu $ of our observer (reference frame), what is the energy of the photon $p^\mu $? We write $p^0$ in the observer rest frame in a coordinate invariant way. We do this by general covariance, so first we write it in locally inertial coordinates. We have the 4-momentum of the observer and photon. We can take the inner product
\[
	-U_{\mu } p^\mu =-U^\mu p^\nu g_{\mu \nu } 
.\]
In the observer's rest frame, $U^\mu =(1,0,0,0)$, so we only get a contribution when $\mu =0$:
\[
	-p^\nu g_{0 \nu } =p^0
.\]
This is the energy of the photon, and is manifestly coordinate invariant since $E=-U_\mu p^\mu $.

Now we solve the geodesic equation. Since massless particles travel along right rays, $\mathrm{d} s^2=0$. We then have
\[
	0=-\left( \frac{\mathrm{d}t}{\mathrm{d}\lambda }  \right) ^2+a^2(t))\left( \frac{\mathrm{d}x}{\mathrm{d}\lambda }  \right) ^2
.\]
Whenever we want to solve for the geodesic equation of a massless particle, start here, since like its kinda like the Geodesic equation. We assume the photon is traveling in the $x$ direction only in the $\lambda $ reference frame by rotational invariance. This is why there is only a contribution from $x$. Solving this equation gives
\[
	\frac{\mathrm{d}t}{\mathrm{d}\lambda } =a(t)\frac{\mathrm{d}x}{\mathrm{d}\lambda } \iff \frac{\mathrm{d}x}{\mathrm{d}\lambda } =\frac{1}{a(t)} \frac{\mathrm{d}t}{\mathrm{d}\lambda } 
.\]
We see that geodesic equation is of the form
\[
	0=t''+a \dot{a} (x')^2=t''+\frac{\dot{a} }{a} (t')^2
.\]
This is solved by $t'=\frac{\omega _0}{a} $ for $\omega _0$ some integration constant.

Recall that the energy of the photon was
\[
	E=-U_\mu p^\mu =-U_\mu (x')^\mu 
.\]
We choose coordinates $U^\mu =(1,0,0,0)$ because the Earth is not moving at the speed of light, and because we can take the entire universe to be moving with respect to us (we are moving slowly through it). So this becomes
\[
	E=g_{00} U^0 t'=\frac{\omega _0}{a} 
.\]
In particular, the energy of a photon \emph{decreases} as the universe expands. This makes sense since the energy is inversely proporitional to the photon, and as it expands, the wavelength expands. The redshift is just
\[
	Z=\frac{\Delta \lambda }{\lambda _i} =\frac{\lambda _f-\lambda _i}{\lambda _i} =\frac{a(t_f)}{a(t_i)} -1
.\]
We choose $a(t)$ for $t=$ today to be 1, so we simply have
\[
	Z=\frac{1}{a(t_i)} -1
.\]

In reality, the universe has things in it. They all have the same form for their stress tensor, which is the perfect fluid stress tensor
\[
	T_{\mu \nu } =\left( p+\rho  \right) U_\mu U_\nu +pg_{\mu \nu } 
.\]
Here, $\rho $ is the energy density and $p$ is the pressure. We can find $p$ given $\rho $ for matter, photons, and dark energy, which are the three things in the universe.

For a single particle $x^\mu (\lambda )$, apparently
\[
	T^{\mu \nu } =\int p^\mu \frac{\mathrm{d}x^\nu }{\mathrm{d}\lambda } \delta ^3(x-x(\lambda )) \,\mathrm{d} \lambda 
.\]
If we choose the rest frame of the particle, then
\[
	T^{00} =p^0
.\]
The formula is manifestly covariant, so we now compute the pressure. Apparently $T^i_j=p\delta ^i_j$. Apparently the answer depends on the temperature of the matter, but in space that is basically 0. For cold matter, $p=0$. So the equation of state
\[
	w=\frac{p}{\rho } =0
\]
for cold matter. For photons, $w=1/3$, so $\rho =3p$. This follows because $p^2$ (momentum) is 0 for massless particles, since the square of the 4-momentum is the rest mass.

For dark energy, we assume dark energy is a cosmological constant, meaning $T^{\mu \nu } $ is proportional to $-\eta ^{\mu \nu } $. This is the only thing we can have without breaking Lorentz invariance. We then see that $w=-1$.

So now, how does $\rho $ change with the expansion of the universe? Intuitively, $\rho \sim 1/a^3$ for matter, since volume is proportional to length cubed. For photons, the wavelength also gets redshifted, so we expect it to be $\rho \sim 1/a^4$. The cosmological constant is... constant... $T^{\mu \nu }=\Lambda \eta ^{\mu \nu } $ and $\Lambda $ is independent.

Without general relativity, $\partial _\mu T^{\mu \nu } =0$. This isn't a tensor, so we have to take the covariant derivative to make it a tensor:
\[
	\nabla _\mu T^{\mu \nu } =0
.\]
This can be solved two different ways. It is easier to solve it if we lower $\nu $:
\[
	\nabla _\mu T^\mu _\nu 
.\]
Of course, we can always relate these by $g_{\mu \nu } $. We have
\[
	\nabla _\mu T^\mu _\nu =\partial _\mu T^\mu _\nu +\Gamma _{\mu \alpha } ^\nu T_\nu ^\alpha -\Gamma _{\mu \nu } ^\alpha T^\mu _\alpha 
.\]
We found peviously that in the FRW metric,
\[
	\Gamma ^t_{ij} =a \dot{a} \delta _{ij} ,\qquad \Gamma _{tj} ^i=\frac{\dot{a} }{a} \delta _j^i
,\]
and all others are 0. Plugging this in,
\[
	0=\partial _0T^0+\partial _iT^i_\nu +\Gamma^0_{0\alpha } T^\alpha _0 +\Gamma_{i\alpha } ^iT^\alpha _0 -\Gamma ^0_{\mu 0} T^\mu _0-?
.\]
Apparently most vanish because of some magic by seeing that off-diagonals vanish on the Christoffel symbols. The nonvanishing terms are
\[
	=\partial _tT^i_t+\Gamma _{it} ^iT_t^t-\Gamma^i_{jt} T_i^j=
.\]
We use $T^\mu _\nu =\operatorname{diag} (\rho ,p,p,p)$ and our definition of the Christoffel symbols to see that this is
\[
	-\dot{\rho } -\frac{\dot{a} }{a} \delta _i^i\rho -\frac{\dot{a} }{a} \delta _j^iw\rho \delta _i^j=-\left( \dot{\rho } +3\frac{\dot{a} }{a} (1+w)\rho  \right) 
.\]
This is solved by
\[
	\rho (t)=\frac{\rho _0}{(a(t))^{3(1+w)} } 
.\]

